\section{Variational inevitability of the golden branch: from minimal discrepancy to canonical clocks}
\label{sec:golden_ratio}

The distinguished role of the golden ratio in HPA--$\Omega$ follows from two independent mechanisms:
\begin{enumerate}
  \item \textbf{Discrepancy minimization.} At finite readout depth, the golden branch minimizes monotone discrepancy proxies controlled by continued-fraction coefficients (Section~\ref{subsec:golden_optimality}). This reduces mismatch accumulation and therefore the total cost required to sustain self-consistent readout.
  \item \textbf{Shortest canonical coding.} In the golden branch, Ostrowski numeration degenerates to Zeckendorf decomposition: every integer has a unique representation as a sum of nonconsecutive Fibonacci numbers \cite{Zeckendorf1972}. This yields a canonical discrete tick structure that naturally decomposes finite budgets across scales without external tuning.
\end{enumerate}

\noindent These two mechanisms are logically independent: the first is about \emph{anti-resonance} (hardness of rational approximation) and mismatch suppression, while the second is about \emph{addressability} (canonical normal form for finite tick budgets). Their conjunction explains why the golden branch repeatedly appears as a stable fixed point across the HPA--$\Omega$ chain: it simultaneously lowers mismatch cost and shortens the bookkeeping needed to maintain readout coherence.

In CAP terms, the golden branch is not an aesthetic ornament; it is a variational attractor selected by readout sustainability under finite information and Weyl complementarity.


